\section{The binary incidence construction and its smooth specializations}
\label{sec:binary-specializations}
The fixed-prime coherent-line theorem strengthens several binary constructions. We give their equations and fields to distinguish what the incidence identity already supplies from what the later smoothness step supplies. Their records, together with the present proof, are indexed in the workbench.

Use the preprocessed width-three formula and incidence construction of Lemma~\ref{lem:incidence}, writing $M=|\mathcal I|$ and renaming $w_{e,b}$ as $z_{e,b}$ for $b\in\{0,1\}$. Thus $P=\prod_fP_f\prod_iQ_i$, $Z=\prod_ez_{e,0}z_{e,1}$, and the lemma gives the integral identity
\begin{equation}\label{eq:ff-legacy-incidence}
 [Z]P=S.
\end{equation}
Every clause factor has at most eight terms. The polynomial $P$ is separately homogeneous of degree two in each pair: the clause factor and the variable factor at that incidence each contribute degree one. Its total degree is $2M$.

Over $\mathbb F_2$ put $r=2M+1$ and $F=z_0P$. This is a degree-$r$ form in $r$ variables, whose squarefree coefficient is $c=S\bmod2$. The projective hypersurface $V(F)$ has $H^0(\mathcal O)=\mathbb F_2$ and one further coherent group, a line in degree $r-2$ with Frobenius scalar $c$, by the hypersurface exact sequence. It may be singular or reducible. Empty clauses and zero-variable formulas have directly known counts and can be replaced by fixed positive-unit or contradictory-unit instances before forming the incidence variables. All ordinary instances have $M$ linear in the formula length and a polynomial-size factored equation.

\subsection{An affine double cover and the first lost binary digit}
There is a related smooth affine construction with quadratic equations. For a three-CNF formula on $n\geq1$ bits with $m$ clauses, let $v_C(x)$ be the product of the false-literal indicators of clause $C$, and put
\[
 A_\Phi(x)=\prod_C(1+v_C(x)),\qquad
 f_\Phi(x)=1+A_\Phi(x)+x_1^{6m+3}+x_1
 \quad\text{over }\mathbb F_2.
\]
On Boolean inputs, $f_\Phi=0$ exactly at satisfying assignments. Its degree is the odd integer $6m+3$, and its leading term is $x_1^{6m+3}$. Hence
\[
 U_\Phi=\{t^2+t=f_\Phi(x)\}\longrightarrow\mathbb A^n_{\mathbb F_2}
\]
is finite \'etale of degree two: its coordinate algebra has basis $1,t$ over the base and derivative one in $t$. It is geometrically integral. Indeed, reducibility over the algebraic closure of the rational function field would give $f_\Phi=h^2+h$ for a rational function $h$. Writing $h=A/B$ in relatively prime polynomials forces $B$ to be constant, by reducing $A^2+AB=f_\Phi B^2$ modulo each prime divisor of $B$. The total degree of $h^2+h$ is then even or zero, a contradiction.

Each satisfying Boolean input has two rational points above it and each nonsatisfying input has none, so $\#U_\Phi(\mathbb F_2)=2S$. A repeated-squaring circuit for $f_\Phi$ has $O(n+m+\log(m+2))$ gates. Giving each gate its own coordinate and graph equation produces a complete intersection of sparse quadratic equations; the Jacobian minor in the gate coordinates followed by $t$ is lower triangular with diagonal one. Successive elimination recovers the original hypersurface, so this presentation preserves the entire scheme. It gives a smooth geometrically integral quadratic complete intersection with an explicit global smoothness certificate. Its count modulo two is always zero, whereas its count modulo four determines $S\bmod2$. Exact counting is hard under division by two, which is a numerical reduction rather than a parsimonious one. No projective compactification with polynomially many ordinary coordinates is asserted for this affine presentation.

\subsection{A pencil, its bad parameters, and a fixed quadratic field}
Put
\[
 H=z_0^r+\sum_{e=1}^{M}(z_{e,0}^r+z_{e,1}^r).
\]
The pencil $V(H+uF)$ is smooth over $\mathbb F_2(u)$, since its fiber at zero is the smooth odd-degree Fermat hypersurface. More explicitly, set $J=r(r-2)+1$ and $T=\binom{r(r-1)}{r-1}$. Every monomial of degree $J$ is divisible by a coordinate to power $r-1$. Multiply the corresponding partial derivative of $H+uF$ by the quotient monomial. The resulting $T$ vectors in degree $J$ form a matrix $I+uA$. Its determinant $\Delta(u)$ has $\Delta(0)=1$ and degree at most $T$. If $\Delta(a)\ne0$, the gradient ideal contains every degree-$J$ monomial and has no projective zero. This proves geometric smoothness without calculating the exponentially large matrix.

The integer $T$ has $O(r\log r)$ bits. A power-of-two extension degree $b=O(r\log r)$ can be chosen with $q=2^b\geq8(T+1)$. The quadratic-tower construction used in Section~\ref{sec:ff} produces the field deterministically. Sample $a$ from an affine half-space avoiding zero. At most $T$ of the $q/2$ possible parameters are roots of $\Delta$, so smoothness has probability greater than $3/4$. This is a randomized finite-field specialization. On its smooth outputs the absolute Frobenius coefficient is $ac$, whereas the $q$-power scalar is $(ac)^{q-1}=c$. Over the generic field the coefficient is $uc$ and the action is semilinear; it is not an endomorphism linear over $\mathbb F_2(u)$.

The incidence structure gives a deterministic specialization with a stronger field control.
\begin{lemma}\label{lem:ff-bad-field}
Let $P$ be a nonzero polynomial over $\mathbb F_2$, separately homogeneous of degree two in each of $M$ variable pairs, and let $r=2M+1$. Every parameter $a$ for which $V(H+a z_0P)$ is singular belongs to
\[
 \mathbb F_{2^{r e}},\qquad e=\operatorname{ord}_r(2).
\]
\end{lemma}
\begin{proof}
At a singular point $z_0$ cannot be zero: the remaining derivative equations would force every other coordinate to vanish. Normalize $z_0=1$. The $z_0$ derivative gives $aP=1$. Euler's identity in each pair, whose degree is two and hence zero in the field, combines the two derivative equations into
$z_{j,0}^r=z_{j,1}^r$. Both coordinates are nonzero, because vanishing of a whole pair would force $P=0$. Write them as $(s_j,s_jt_j)$. Then $t_j^r=1$, so all $t_j$ belong to $K=\mathbb F_{2^e}$.

Pairwise homogeneity gives $P=(\prod_j s_j^2)C(t)$. It likewise writes $z_{j,0}P_{z_{j,0}}=(\prod_i s_i^2)C_j(t)$. Multiplication of the corresponding derivative equation by $s_j$, followed by $aP=1$, gives
\[
 s_j^r=C_j(t)/C(t)\in K^*.
\]
If $q_0=2^e$, then $r(q_0-1)$ divides $q_0^r-1$, because $q_0\equiv1\pmod r$. Thus each $s_j$ belongs to $\mathbb F_{q_0^r}$. All normalized coordinates, and then $a=1/P$, belong to this field.
\end{proof}

Compute $e$ by repeated doubling modulo $r$, choose the least power of two $b>re$, and use the last generator $\theta$ of a quadratic tower of degree $b$. It has absolute degree $b$, whereas every bad parameter has degree at most $re$. Hence $V(H+\theta F)$ is smooth. Here $b\leq2re<2r^2$, the field presentation is polynomial in size, and its top coherent line has absolute coefficient $\theta c$ and $2^b$-power scalar $c$.

There is also a specialization over the one fixed field $K=\mathbb F_4=\mathbb F_2(\omega)$, with $\omega^2+\omega+1=0$. Pad the formula by forced-one variables until $r=2^k-1$ for an odd $k\geq3$. The least such $r$ is less than $4(r_0+1)$ for the original degree $r_0$, and each new variable adds exactly one incidence and has a unique legal value. Thus the padding is linear and preserves $S$. Now $e=k$, so the bad-parameter field has odd degree $rk$. It contains neither $\omega$ nor $\omega^2$. Consequently
\[
 X=V(G),\qquad G=H+\omega F,
 \qquad X\subset\mathbb P_K^{r-1},
\]
is smooth and geometrically integral. Its coherent groups are $K$ in degrees zero and $d=r-2$, and zero elsewhere. Absolute Frobenius has coefficient $\omega c$ on the latter line, and $4$-power Frobenius has scalar $c$. These are distinct assertions about the same semilinear map. The point-count parity is $1+c$: the coefficient of the cube of the coordinate product in $G^3=GG^2$ is $(\omega c)^3=c$, since the term from $G$ must have positive odd exponent in every coordinate and total degree $r$, hence be squarefree. Summing $1-G^3$ on the affine cone over $\mathbb F_4$ then gives the projective congruence.

\subsection{Descent to a four-dimensional coherent realization over the prime field}
One can descend the preceding construction to a smooth projective scheme over $\mathbb F_2$ with total coherent dimension four. Use $r^2$ projective coordinates $a_i,b_{ij},c_{ij}$, $i<j$, and form the Hermitian matrix
\[
 M_{ii}=a_i,\qquad M_{ij}=b_{ij}+\omega c_{ij},
 \qquad M_{ji}=b_{ij}+\omega^2c_{ij}.
\]
Define $W\subset\mathbb P^{r^2-1}_{\mathbb F_2}$ by both binary coefficient components of every two-by-two minor of $M$ and of each column polynomial
\[
 G(M_{1j},\ldots,M_{rj}),\qquad 1\leq j\leq r.
\]
There are $2\binom r2^2$ quadratic equations and $2r$ equations of degree $r$. An arithmetic operation over $K$ becomes a constant number of operations on its two binary coordinates, so this is a polynomial-size homogeneous circuit presentation.

Over $K$ the entries of $M$ are independent coordinates: the change from $(b_{ij},c_{ij})$ to $(M_{ij},M_{ji})$ has determinant $\omega^2-\omega=1$. Imposing both components of a polynomial is equivalent to imposing it and its coefficientwise conjugate. The minors therefore define exactly the rank-one projective matrix scheme. To check scheme structure, on a chart with $M_{ij}=1$ they give $M_{kl}=M_{kj}M_{il}$, leaving precisely the free coordinates of the corresponding chart of $\mathbb P^{r-1}\times\mathbb P^{r-1}$. These charts cover, including over nonreduced test algebras.

Writing $M=zw^{\mathsf t}$, the column equations become $w_j^rG(z)=0$, and their conjugates become $z_i^rG^\sigma(w)=0$. On every product chart some $z_i$ and $w_j$ are units, so the generated ideal sheaf is exactly $(G(z),G^\sigma(w))$. Thus
\[
 W_K\simeq X\times_K X^\sigma.
\]
Coefficient conjugation transposes $M$ and exchanges the two factors. This is the descent datum of $\operatorname{Res}_{K/\mathbb F_2}X$, since $K\otimes_{\mathbb F_2}K=K\times K$. Faithfully flat descent identifies $W$ with this restriction of scalars. Smoothness descends from the product. Its geometric global functions form a field, by the coherent calculation and the tensor product of affine Cech complexes, so it is geometrically connected. A smooth connected scheme over an algebraically closed field is integral: its regular local rings are domains, which prevents distinct irreducible components from intersecting. Hence $W$ is geometrically integral.

Let $e,e'$ be the top coherent generators on the two factors. The same tensor product of Cech complexes shows that the only groups of $W_K$ are $K$, $Ke\oplus Ke'$, and $K(e\cup e')$, in degrees $0,d,2d$. The descended bases are
\[
 1,\qquad v_1=e+e',\quad v_2=\omega e+\omega^2e',
 \qquad w=e\cup e'.
\]
They are fixed by conjugation and factor exchange, and their coefficient determinant is one. Absolute Frobenius is semilinear over $K$, with
$F(e)=\omega c e$ and $F(e')=\omega^2c e'$. Therefore
\[
 F(v_1)=cv_2,\qquad F(v_2)=cv_1,\qquad F(w)=cw.
\]
The coherent dimensions over $\mathbb F_2$ are exactly $1,2,1$, and the Frobenius matrices are $1$, $c\left(\begin{smallmatrix}0&1\\1&0\end{smallmatrix}\right)$, and $c$. In particular the top coherent scalar is still the parity count.

Finally $W(\mathbb F_2)$ is in bijection with $X(\mathbb F_4)$. The forward map sends $[z]$ to $z\sigma(z)^{\mathsf t}$; scaling $z$ multiplies this matrix by its scalar norm, which is one in $\mathbb F_2^*$. Conversely a nonzero binary Hermitian rank-one matrix has a nonzero diagonal entry, since otherwise its minors force every entry to vanish. A column through a diagonal entry one reconstructs $z$, and the column equation gives $G(z)=0$. Thus the rational-point parity is $1+c$ as well. The main fixed-prime hypersurface theorem achieves the same scalar on one coherent line directly over the prime field; the construction here explains the earlier four-dimensional realization and its different input equations.
